\Article{3}
Nehmen wir anstelle der diskreten Punkte $M,M',M''$ usw. einen kontinuierlichen Körper an, der einen Raumbereich $S$ mit gleichmäßiger Dichte $=C$ ausfüllt, so geht die Summe
\[
M\Phi(m,M)+M'\Phi(m,M')+M''\Phi(m,M'')+\ \text{etc.}
\]
in das Integral $C\int\mathrm{d}S.\Phi(m,\mathrm{d}S)$ über, das sich über den ganzen Raumbereich $S$ erstreckt. Entsprechend bezeichnen wir mit $(m,\mathrm{d}S)$ den Abstand des Punkts $m$ von einem beliebigen Element des Raumbereichs $S$, das mit $\mathrm{d}S$ bezeichnet wird.

Betrachten wir darüber hinaus anstelle der diskreten Punkte $m,m',m''$ usw. einen kontinuierlichen Körper, der einen Raumbereich $s$ mit gleichmäßiger Dichte $=c$ ausfüllt, so erfordert die Berechnung von $\Omega$ eine doppelte Integration. Sie ist so auszuführen, dass zunächst für einen unbestimmten Punkt $\mu$ der Wert des Ausdrucks
\SourceLeaf{42}
\[
-gz+\tfrac12c\int\mathrm{d}s.\varphi(\mu,\mathrm{d}s)+C\int\mathrm{d}S.\Phi(\mu,\mathrm{d}S)
\]
ermittelt wird. Dabei ist $z$ die Höhe des Punkts $\mu$ über der Ebene $H$. Das erste Integral ist über den ganzen Raumbereich $s$, das zweite über den ganzen Raumbereich $S$ zu erstrecken. Dieser Wert hängt allein von der Lage des Punkts $\mu$ ab. Bezeichnet man ihn mit $[\mu]$, so ergibt sich
\[
\Omega=c\int\mathrm{d}s.[\mathrm{d}s]
\]
wobei sich die Integration über den ganzen Raumbereich $s$ erstreckt.

Kürzer lässt sich das so ausdrücken:
\[
\Omega=-gc\int z\mathrm{d}s+\tfrac12cc\int\int\mathrm{d}s.\mathrm{d}s'.\varphi(\mathrm{d}s,\mathrm{d}s')+cC\int\int\mathrm{d}s.\mathrm{d}S.\Phi(\mathrm{d}s,\mathrm{d}S)
\]
Dabei bezeichnen $s,s'$ eigentlich ein und denselben vom beweglichen Körper ausgefüllten Raumbereich, der jedoch für die doppelte Integration zweimal in seine Elemente zerlegt werden muss.

\Article{4}
Das Kennzeichen fluider Körper ist die vollkommene Beweglichkeit selbst ihrer kleinsten Teile, sodass sie jede Gestalt annehmen können und selbst der geringsten Kraft nachgeben, die ihre Gestalt zu verändern sucht. Bei nicht ausdehnbaren Fluiden, also Flüssigkeiten, denen unsere Untersuchung gewidmet ist, muss das Volumen jedes Teilchens bei allen Gestaltänderungen konstant bleiben. Betrachten wir daher einen fluiden Körper, dessen Bewegung durch einen unbeweglichen festen Körper, ein Gefäß, begrenzt wird. Auf seine Teilchen wirken unserer Annahme nach außer der Schwerkraft sowohl die gegenseitige Anziehung der Teilchen als auch die Anziehung der Gefäßteile. Der Gleichgewichtszustand erfordert, dass der Wert von $\Omega$ ein Maximum ist, das heißt, dass keine unendlich kleine Umlagerung der Flüssigkeitsteile $\Omega$ einen positiven Zuwachs geben kann. Der Wert von $\Omega$ kann sich offenbar nur ändern, soweit sich die Gestalt des von der ganzen Flüssigkeit ausgefüllten Raumbereichs ändert. Eine Bewegung allein im Inneren der Flüssigkeit kann ihn nicht ändern. Gleichgewicht liegt demnach vor, wenn $\Omega$ bei keiner unendlich kleinen, mit der Gestalt des Gefäßes vereinbaren Änderung jener Gestalt bei konstant bleibendem Volumen zunehmen kann. Daraus folgt unmittelbar, dass jene auf die Flüssigkeit wirkenden Kräfte keine innere Bewegung der Flüssigkeit hervorrufen können und sich gegenseitig im Gleichgewicht halten, wenn die Gestalt sich überhaupt nicht ändern kann, weil das Gefäß die Flüssigkeit von allen Seiten umschließt und berührt.

\Article{5}
Wir gehen zur genaueren Untersuchung des Ausdrucks $\Omega$ über, der als Grundlage der Theorie des Gleichgewichts von Flüssigkeiten anzusehen ist. Beginnen wir \SourceLeaf{43}mit dem ersten Term, so ist unmittelbar klar, dass $\int z\mathrm{d}s$ das Produkt aus dem Volumen des Raumbereichs $s$ und der Höhe seines Schwerpunkts über der Ebene $H$ darstellt. Die Produkte derselben Höhe mit der Masse beziehungsweise mit dem Gewicht der Flüssigkeit sind demnach $c\int z\mathrm{d}s$ und $gc\int z\mathrm{d}s$. Wären also die Teile der Flüssigkeit außer der Schwerkraft keiner weiteren Kraft ausgesetzt, so müsste die Höhe des Schwerpunkts im Gleichgewicht so klein wie möglich sein. Daraus lässt sich leicht schließen, dass der freie Teil der Oberfläche oder ihre freien Teile in ein und derselben horizontalen Ebene liegen müssen, welche die Flüssigkeit von oben begrenzt.

\Article{6}
Die Entwicklung des zweiten und dritten Terms betrifft zwei Sonderfälle einer allgemeinen Aufgabe. Sind zwei beliebige Raumbereiche gegeben, so ist jedes Element des ersten mit jedem Element des zweiten zu paaren, und die Produkte aus je drei Faktoren sind zu summieren. Die Faktoren sind das Volumen des Elements aus dem ersten Raumbereich, das Volumen des Elements aus dem zweiten und eine gegebene Funktion ihres gegenseitigen Abstands. Der zweite Term betrifft den Fall, in dem beide Raumbereiche identisch sind, der dritte den Fall, in dem der eine ganz außerhalb des anderen liegt. Die vollständige Aufgabe umfasst zwei weitere Fälle, in denen der eine Raumbereich ein Teil des anderen ist oder beide einen Teil gemeinsam haben. Die ersten beiden Fälle könnten zwar für unser Vorhaben genügen und die beiden übrigen ließen sich leicht auf sie zurückführen. Dennoch lohnt es sich, die schon für sich genommen bemerkenswerte Aufgabe in voller Allgemeinheit zu behandeln. Bei dieser allgemeinen Untersuchung bezeichnen wir die Raumbereiche mit $s,S$ und die Funktion des Abstands mit dem Zeichen $\varphi$. Bei der Anwendung auf den zweiten Term ist anstelle von $S$ der Raumbereich $s$ selbst und bei der Anwendung auf den dritten Term anstelle der Funktion $\varphi$ die Funktion $\Phi$ selbst einzusetzen. Es geht also um das Integral
\[
\int\int\mathrm{d}s.\mathrm{d}S.\varphi(\mathrm{d}s,\mathrm{d}S)
\]
das wie eine doppelte Integration aussieht. Tatsächlich hängen die Elemente jedes der beiden Raumbereiche von drei Variablen ab, sodass es eine sechsfache Integration umfasst, die wir nun auf eine vierfache zurückführen werden.

